\section{Galactic rotation curves: TGP vs CDM and MOND}
\label{sec:rotation_curve}

\subsection{Setup}

We compute rotation curves for a Milky-Way-like galaxy under four
scenarios: baryons only (Newtonian), TGP with Yukawa screening,
CDM with NFW halo, and MOND.

\paragraph{Galaxy model.}
Exponential stellar disk ($M_{\rm disk}=5\times10^{10}\,M_\odot$,
$R_d=3.5\,{\rm kpc}$), Hernquist bulge
($M_{\rm bulge}=10^{10}\,M_\odot$, $r_b=0.7\,{\rm kpc}$),
gas disk ($M_{\rm gas}=5\times10^9\,M_\odot$, $R_{\rm gas}=7\,{\rm kpc}$).
NFW halo: $M_{200}=10^{12}\,M_\odot$, $c_{\rm NFW}=15$,
$r_s=13.6\,{\rm kpc}$.
MOND acceleration constant $a_0=1.2\times10^{-10}\,{\rm m\,s}^{-2}$.

\paragraph{TGP rotation curve.}
For TGP Path~B, the 2-body coupling is $C=m/(2\sqrt{\pi}m_{\rm Pl})$,
giving $C^2=Gm_1m_2$ — identical to Newtonian gravity.
The radial acceleration from TGP is
\begin{equation}
  a_{\rm TGP}(r) = \frac{GM(r)}{r^2}\,e^{-m_{\rm sp}r}(1+m_{\rm sp}r),
\end{equation}
where the Yukawa screening factor $f(r)=e^{-m_{\rm sp}r}(1+m_{\rm sp}r)$
satisfies $f(0)=1$ and $f(r)\to0$ exponentially for $r\gg1/m_{\rm sp}$.

\subsection{Results}

\begin{table}[h]
\centering
\caption{Circular velocity $v_{\rm circ}$ [km/s] for each model.
TGP with any $m_{\rm sp}>0$ gives \emph{less} gravity than Newton at
large $r$ — opposite to what flat rotation curves require.}
\label{tab:rotation_curve}
\begin{tabular}{rllllll}
\hline
$r$ & Baryons & TGP & TGP & TGP & CDM & MOND \\
[kpc] & (Newton) & $\lambda=1$ & $\lambda=5$ & $\lambda=15$ & +NFW & \\
& & [kpc] & [kpc] & [kpc] & & \\
\hline
$ 1$ & $138$ & $118$ & $137$ & $138$ & $157$ & $139$ \\
$ 5$ & $170$ & $ 34$ & $145$ & $166$ & $222$ & $182$ \\
$10$ & $165$ & $  4$ & $105$ & $153$ & $239$ & $195$ \\
$15$ & $145$ & $  0$ & $ 65$ & $124$ & $235$ & $194$ \\
$20$ & $126$ & $  0$ & $ 38$ & $ 99$ & $228$ & $190$ \\
$30$ & $100$ & $  0$ & $ 13$ & $ 64$ & $217$ & $185$ \\
\hline
\end{tabular}
\end{table}

\begin{table}[h]
\centering
\caption{Yukawa screening factor $f(r)=e^{-m_{\rm sp}r}(1+m_{\rm sp}r)$.
At $r=20\,{\rm kpc}$, even $\lambda=15\,{\rm kpc}$ gives $f=0.62$ —
a $38\%$ reduction in gravity relative to Newton.}
\label{tab:yukawa_screening}
\begin{tabular}{rrrrr}
\hline
$r$ [kpc] & $\lambda=1\,{\rm kpc}$ & $\lambda=5\,{\rm kpc}$ &
$\lambda=15\,{\rm kpc}$ \\
\hline
$ 1$ & $0.736$ & $0.983$ & $0.998$ \\
$ 5$ & $0.040$ & $0.736$ & $0.955$ \\
$10$ & $0.0005$ & $0.406$ & $0.856$ \\
$20$ & $\approx0$ & $0.092$ & $0.615$ \\
$30$ & $\approx0$ & $0.017$ & $0.406$ \\
\hline
\end{tabular}
\end{table}

\subsection{Analysis}

\paragraph{Why TGP cannot explain flat rotation curves.}
The Yukawa potential $e^{-m_{\rm sp}r}/r$ satisfies
\begin{equation}
  \frac{e^{-m_{\rm sp}r}}{r} < \frac{1}{r} \quad \forall\,m_{\rm sp}>0,\;r>0.
\end{equation}
Therefore TGP acceleration is \emph{always less than or equal to}
Newtonian.  Flat rotation curves ($v={\rm const}$) require
$a(r)\propto1/r$ — more gravity at large $r$ than the baryonic
Newtonian prediction.  TGP provides systematically less gravity,
moving in the \emph{wrong direction}.

\paragraph{CDM (NFW halo) works.}
The NFW halo provides an additional $\sim100\,{\rm km/s}$ at $r=20\,{\rm kpc}$,
sustaining a flat curve of $v\approx220\,{\rm km/s}$ consistent with
the observed Milky Way rotation velocity.

\paragraph{MOND works.}
In the deep-MOND limit ($a\ll a_0$), Milgrom's relation gives
$a^2/a_0=a_N\Rightarrow v^4=GMa_0$ (Tully-Fisher relation).
For $M_{\rm bar}=6.5\times10^{10}\,M_\odot$, the MOND prediction is
$v_{\rm TF}=(GM_{\rm bar}a_0)^{1/4}=179\,{\rm km/s}$, consistent with
the computed curve $v(30\,{\rm kpc})=185\,{\rm km/s}$.

\paragraph{Three-body TGP correction.}
For ordinary baryonic matter ($C\sim10^{-19}$), the three-body
coupling is $C^3\sim10^{-57}$, completely negligible for galactic dynamics.
Even summing over $N\sim10^{11}$ stellar triples within a 1~kpc sphere
gives a relative correction $\delta a_{\rm 3B}/a\sim NC^3\sim10^{-46}$.

\subsection{Conclusion}

TGP Path~B (for baryonic matter with $C\sim10^{-19}$) is
indistinguishable from screened Newtonian gravity.  It:
\begin{enumerate}
\item Reproduces Newtonian gravity exactly at $r\ll\lambda_{\rm sp}$.
\item Predicts \emph{less} gravity than Newton at $r\gtrsim\lambda_{\rm sp}$.
\item Cannot produce flat rotation curves for any $m_{\rm sp}>0$.
\item Has negligible ($\sim10^{-57}$) three-body corrections for baryons.
\end{enumerate}

TGP cannot substitute for dark matter or MOND in explaining rotation curves.
A TGP dark-matter candidate would require near-Planck-mass particles
($C\sim1$, $m\sim m_{\rm Pl}$), but their interaction with baryons
(coupling $\propto C_{\rm baryon}\cdot C_{\rm DM}\sim10^{-19}$) would
be observationally indistinguishable from zero.

The Yukawa screening is a generic property: any Lagrangian with a
\emph{positive} scalar mass $m_{\rm sp}^2>0$ yields a screened potential.
Only a tachyonic scalar ($m_{\rm sp}^2<0$) or a massless scalar ($m_{\rm sp}=0$)
could produce long-range enhancement — but these lie outside TGP Path~B.
